\documentclass[10pt]{article}

\usepackage{amsmath,amssymb,amsthm}
\usepackage{url}

\newtheorem{definition}{Definition}
\newtheorem{problem}{Problem}
\newtheorem{theorem}{Theorem}
\newtheorem{lemma}[theorem]{Lemma}
\newtheorem{remark}[theorem]{Remark}

\newcommand{\Kuhn}{\text{Kuhn}}
\newcommand{\BKuhn}{\text{BlockingKuhn}}
\newcommand{\Ex}{\mathbb{E}}
\newcommand{\ONCE}{\textup{\textsc{once}}}
\newcommand{\EACH}{\textup{\textsc{each}}}

\title{A Lower Bound for Randomized Kuhn's Algorithm on Sparse Graphs - Hard}
\author{Nikita Gaevoy}
\date{}

\begin{document}

\maketitle

\section{Preliminaries}

\subsection{Matchings and augmenting paths}

Let $G=(\mathcal{L}\sqcup\mathcal{R},\,\mathcal{E})$ be a bipartite graph. We
write $L=|\mathcal{L}|$, $R=|\mathcal{R}|$, $V=L+R$ and $E=|\mathcal{E}|$, so
that $L$, $R$, $V$, $E$ are \emph{cardinalities}; the vertices of $\mathcal{L}$
are numbered by $[0,L)$ and those of $\mathcal{R}$ by $[0,R)$. The graph is
assumed connected, so $E\ge V-1$, and it is given by adjacency lists
$\mathrm{neigh}[v]$, $v\in\mathcal{L}$, which list the neighbours of $v$ in
$\mathcal{R}$ in some order. That order is part of the input, and it matters:
the algorithms below read the lists in the order given.

A \emph{matching} is a set $M\subseteq\mathcal{E}$ of pairwise disjoint edges.
A vertex is \emph{free} if no edge of $M$ meets it. An \emph{augmenting path}
for $M$ is a path whose endpoints are both free and whose edges alternate
between $\mathcal{E}\setminus M$ and $M$; replacing $M$ by $M\triangle P$
enlarges the matching by one. By Berge's theorem~\cite{Berge57}, $M$ is maximum
if and only if it admits no augmenting path.

\subsection{Kuhn's algorithm}

\begin{definition}[The search]\label{def:dfs}
Let $\mathit{lmatch}$ and $\mathit{rmatch}$ hold the current matching, with
$-1$ for a free vertex, and let $\mathit{used}$ be an array of flags on
$\mathcal{L}$. The search is
\begin{verbatim}
    dfs_Kuhn(v):
        if used[v]: return false
        used[v] = true
        for dest in neigh[v]:
            if rmatch[dest] == -1 or dfs_Kuhn(rmatch[dest]):
                lmatch[v]    = dest
                rmatch[dest] = v
                return true
        return false
\end{verbatim}
It returns true exactly when it has found an augmenting path starting at $v$ and
has flipped it.
\end{definition}

\begin{definition}[$\Kuhn$]\label{def:kuhn}
The classical algorithm scans the left side once, resetting the flags before
each search:
\begin{verbatim}
    for i in [0, L):
        used = [false] * L
        dfs_Kuhn(i)
\end{verbatim}
\end{definition}

\begin{remark}[On the name]\label{rem:name}
The augmenting-path algorithm of Definition~\ref{def:kuhn} is called
\emph{Kuhn's algorithm} in the competitive-programming literature, after the
Hungarian method~\cite{Kuhn55}; it has no standard name in English-language
textbooks, where it appears as the unadorned augmenting-path or
Ford--Fulkerson-style algorithm for bipartite matching. We keep the name
because it is how the problem below is posed in its source~\cite{Kaban19}.
\end{remark}

$\Kuhn$ runs in $O(VE)$ time, and the bound is tight essentially independently
of the relation between $V$ and $E$ --- for sparse graphs, for dense graphs, and
for everything in between~\cite[\S1]{Kaban19}.

\subsection{The blocking, randomized variant}

Two changes make the algorithm dramatically faster in practice. The first,
due to Burunduk1~\cite{Burunduk1} and described in English by
adamant~\cite{Adamant23}, is to reset the flags once per \emph{round} rather
than once per vertex, so that one $O(E)$ round flips a whole set of
vertex-disjoint augmenting paths instead of at most one. The second is to
randomise the order in which the left vertices and their adjacency lists are
read.

\begin{definition}[$\BKuhn$]\label{def:alg}
Fix a shuffle mode $\sigma\in\{\ONCE,\EACH\}$. The algorithm is
\begin{verbatim}
    reshuffle():
        shuffle(neigh); for row in neigh: shuffle(row)

    lmatch = [-1] * L;  rmatch = [-1] * R
    reshuffle()
    while true:
        if sigma == EACH: reshuffle()
        used  = [false] * L
        found = false
        for i in [0, L):
            if lmatch[i] == -1: found |= dfs_Kuhn(i)
        if not found: break
    return L - (number of -1 entries in lmatch)
\end{verbatim}
Here $\mathrm{shuffle}$ applies a uniformly random permutation, independently
of everything else, and one iteration of the \texttt{while} loop is a
\emph{round}. Under $\sigma=\ONCE$ the two shuffles are performed once,
before the first round; under $\sigma=\EACH$ they are repeated at the
start of every round, which is the ``stronger'', more random variant. Shuffling
$\mathrm{neigh}$ renumbers the left vertices, so the loop
\texttt{for i in [0,L)} then visits them in a uniformly random order.
\end{definition}

\begin{remark}[Correctness and the trivial bound]\label{rem:correct}
$\BKuhn$ returns the maximum matching size, for either mode and every outcome
of the shuffles. Indeed, consider the final round: it performs no augmentation,
so throughout it the matching is fixed, and sharing $\mathit{used}$ across the
searches of a single round is then exactly the sharing that is sound for a fixed
matching --- a left vertex from which no augmenting path exists cannot acquire
one while nothing changes. Every free left vertex therefore fails, no augmenting
path exists, and Berge's theorem~\cite{Berge57} applies.

A round costs $O(E)$ and, unless it is the last, augments at least once, so
there are at most $L+1$ rounds and the running time is $O(VE)$ --- the same
bound as for $\Kuhn$.
\end{remark}

\begin{remark}[Why the Hopcroft--Karp bound does not apply]\label{rem:hk}
The set of augmenting paths flipped in one round is vertex-disjoint, so a round
resembles a phase of the Hopcroft--Karp algorithm~\cite{HopcroftKarp73}, which
runs in $O(E\sqrt V)$. It is not one. Hopcroft and Karp flip a maximal set of
\emph{shortest} augmenting paths, and it is shortness that forces the length of
a shortest augmenting path to increase from phase to phase and so caps the
number of phases at $O(\sqrt V)$. The depth-first round of
Definition~\ref{def:alg} flips paths of arbitrary length --- it commits to the
first augmenting path the search stumbles on --- and no bound on the number of
rounds better than the trivial $L+1$ is known. Adding a breadth-first layering
pass in front of the round restores shortness and turns the algorithm into
Hopcroft--Karp~\cite{Adamant23}, at which point $O(E\sqrt V)$ is available
again; but that is a different algorithm, and in practice it is the one the
heuristic is used \emph{instead} of.
\end{remark}

\subsection{The cost measure}

\begin{definition}[Work]\label{def:work}
For a run of $\BKuhn$ let $\rho$ be the number of rounds and let $I$ be the
total number of iterations of the loop \texttt{for dest in neigh[v]} over the
whole run, that is, the number of \emph{edge inspections}. The \emph{work} of
the run is
\[
  W \;=\; \rho L+I .
\]
The two terms are the flag resets and the searches; everything else costs
$O(V+E)$ once, so the running time on a unit-cost RAM is $\Theta(V+E+W)$.
\end{definition}

The shuffles are the only source of randomness, so $W$ is a random variable
depending on $G$, on the adjacency-list order in which $G$ is presented, and on
the mode $\sigma$. We study its expectation, in the worst case over inputs:
\[
  \mathcal{W}_\sigma(V,E)
  \;=\;
  \max_{G}\ \Ex\bigl[W\bigr],
\]
the maximum being over connected bipartite graphs with $V$ vertices and $E$
edges, presented in an arbitrary order. Shuffling makes
$\mathcal{W}_\sigma$ independent of the presentation order for
$\sigma=\ONCE$ as well, since the first shuffle destroys it; the
adversary chooses the graph only.

\section{Problem formulation}

\begin{problem}[Lower bound]\label{prob:lower}
Prove that $\BKuhn$ is not near-linear on sparse graphs. Formally: exhibit an
infinite family $(G_k)_{k\ge1}$ of connected bipartite graphs with
$V_k\to\infty$, $E_k=O(V_k\log^{O(1)}V_k)$, and
\[
  \Ex\bigl[W(G_k)\bigr]=\omega\bigl(E_k\log V_k\bigr)
\]
for at least one of the two shuffle modes $\sigma$.

The bar $E\log V$ is where the known constructions stop --- it is attained on a
path, and it is also what a random graph gives, since there every non-maximum
matching admits an augmenting path of length $O(\log V)$ with high
probability~\cite{BMST06}; the target is $\Omega(E^{3/2})$, which together with
the companion upper bound would settle the sparse case exactly.
\end{problem}

\begin{remark}[The two modes]\label{rem:modes}
A family defeating either mode is a solution. The modes are not comparable, so
it is worth being clear which way a result is the stronger one: a family
defeating $\sigma=\EACH$ is the stronger statement, since it must survive fresh
randomness at the start of every round, and it does not follow from a family
defeating $\sigma=\ONCE$ or imply one. A solution should say which mode it
treats.
\end{remark}

\begin{remark}[Why the problem is posed on sparse graphs]\label{rem:dense}
The dense case is nearly settled and is not what is asked.
For $E=\Theta(V^{2})$ the reference~\cite{Kaban19} reports a family on which
$\BKuhn$ takes $\Omega\bigl(V^{3}/\log V\bigr)=\Omega(VE/\log V)$ time, which is
the trivial upper bound up to a logarithm, so there is nothing left to win
there. The problem therefore lives in the sparse regime, where the gap runs from
$E\log V$ to $E^{2}$ --- a factor of $V/\log V$ --- and where the source reports
that generalising the dense construction did not work~\cite{Kaban19}.
\end{remark}

\begin{remark}[What does not count]\label{rem:notcount}
Two near misses.
\begin{itemize}
  \item A lower bound for the \emph{unshuffled} algorithm. Sparse families
  forcing $\Omega(V^{2})$ without shuffling are known and easy --- a tree built
  from a long path with a path of length $3$ hanging from every vertex of it,
  labelled so that each round finds one long augmenting
  path~\cite{Dacin21}; the whole content of Problem~\ref{prob:lower} is
  that the shuffle must not rescue the construction, so the expectation in the
  statement is over the shuffles and may not be dropped.
  \item An average-case result. On random graphs the behaviour is understood
  and near-linear --- in a random graph every non-maximum matching admits an
  augmenting path of length $O(\log V)$ with high probability, so
  augmenting-path algorithms run in $O(E\log V)$ there~\cite{BMST06}; the problem is worst case over $G$
  and averages only over the algorithm's own coins, so a family drawn at random
  cannot be one.
\end{itemize}
\end{remark}

\begin{remark}[Variants that are equally interesting]\label{rem:variants}
Implementations often shuffle only the adjacency lists, or only the vertex
order, rather than both; and they often start from a greedy matching. A
solution for any such variant, clearly stated, is worth having.
Measured on the one hard construction available --- that of~\cite{Dacin21} ---
either randomisation alone is already enough to defeat it, so it is not known
which of them the bound actually needs, and no implication between the variants
should be assumed in either direction.
\end{remark}

\begin{remark}\label{rem:dichotomy}
A separate problem asks for the bound from
the other side: an implementation-independent $T(V,E)=o(VE)$, with
$O(E\sqrt E\,)$ as the target. The two are not alternatives and not a
dichotomy. On sparse graphs the truth lies somewhere between $E\log V$ and
$E^{2}$, and most of that range satisfies both at once: a proof that
$\Ex[W]=\Theta(E^{3/2})$, the target named in both, would solve this problem and
the companion together. Only a bound of $O(E\log V)$ on the companion side would
exclude this one.
\end{remark}

\begin{thebibliography}{BMST06}

\bibitem[{ada}23]{Adamant23}
{adamant}.
\newblock Implementing {D}initz on bipartite graphs.
\newblock Codeforces blog entry 118098, 2023.

\bibitem[Ber57]{Berge57}
Claude Berge.
\newblock Two theorems in graph theory.
\newblock {\em Proceedings of the National Academy of Sciences},
  43(9):842--844, 1957.

\bibitem[BMST06]{BMST06}
Holger Bast, Kurt Mehlhorn, Guido Sch\"{a}fer, and Hisao Tamaki.
\newblock Matching algorithms are fast in sparse random graphs.
\newblock {\em Theory of Computing Systems}, 39(1):3--14, 2006.

\bibitem[{Bur}15]{Burunduk1}
{Burunduk1}.
\newblock Matching. {F}ast {K}uhn.
\newblock Codeforces blog entry 17023, in Russian, 2015.

\bibitem[{dac}17]{Dacin21}
{dacin21}.
\newblock An {$\Omega(V^2)$} worst case for blocking {K}uhn on a sparse graph.
\newblock Comment 417533 on Codeforces blog entry 58048, 2017.

\bibitem[HK73]{HopcroftKarp73}
John~E. Hopcroft and Richard~M. Karp.
\newblock An {$n^{5/2}$} algorithm for maximum matchings in bipartite graphs.
\newblock {\em SIAM Journal on Computing}, 2(4):225--231, 1973.

\bibitem[{Kab}19]{Kaban19}
{Kaban-5}.
\newblock What is the exact time complexity of randomized {K}uhn's algorithm?
\newblock Computer Science Stack Exchange, question 115411, 2019.
\newblock Version of 2019-10-03.

\bibitem[Kuh55]{Kuhn55}
H.~W. Kuhn.
\newblock The {H}ungarian method for the assignment problem.
\newblock {\em Naval Research Logistics Quarterly}, 2(1--2):83--97, 1955.

\end{thebibliography}

\end{document}
